\documentclass[10pt,a4paper]{amsart}
\usepackage[margin=19mm]{geometry}
\usepackage{amsmath,amssymb,mathtools}
\usepackage[scr=rsfs,cal=euler]{mathalfa}
\usepackage{xfrac}
\usepackage[unicode,hidelinks,bookmarks=false]{hyperref}
\newcommand{\ie}{i.e.\ }
\newcommand{\cf}{cf.\ }
\newcommand{\resp}{respectively\ }
\newcommand{\Subsection}{\subsection}
\providecommand{\nobreakdash}{\nobreak\hbox{-}}
\setlength{\emergencystretch}{2em}
\multlinegap0pt
\usepackage{bm}
\usepackage{bbm}
\usepackage{dsfont}
\usepackage{xspace}
\usepackage{cancel}
\usepackage{mathtools}
\usepackage{amsthm}
\makeatletter
\@ifundefined{equation}{\newtheorem{equation}{Equation}}{}
\@ifundefined{lemma}{\newtheorem{lemma}[equation]{Lemma}}{}
\@ifundefined{remark}{\newtheorem{remark}[equation]{Remark}}{}
\makeatother
\title[On 2-distance 16-coloring of planar graphs with maximum degr]{On 2-distance 16-coloring of planar graphs with maximum degree at most five}
\author{Zakir Deniz}
\date{Source: arXiv:2310.15899v3 (CC BY 4.0)}
\begin{document}
\maketitle

\noindent\emph{Source and licence.} On 2-distance 16-coloring of planar graphs with maximum degree at most five. Version arXiv:2310.15899v3 (CC BY 4.0), distributed under the Creative Commons Attribution 4.0 International licence, as recorded in the arXiv licence metadata for this exact version and shown on the abstract page https://arxiv.org/abs/2310.15899v3. Full rights notice: licenses/arxiv-2310.15899-CC-BY-4.0.txt.\par\smallskip

\noindent\textit{Editorial scope.} Counted unit: the Lemma whose source label is \texttt{lem:support-properties} in the article (source0.tex lines 1063--1071, source label \texttt{lem:support-properties}), delivered together with the article's own complete original proof of it (source0.tex lines 1072--1098, one \texttt{proof} environment of 27 lines). No mathematics is written, repaired or re-derived: statement and proof are the article's own text line for line. Required context retained uncounted: lem:faces-in-3-face (lines 320--322); rem:specil vertices (lines 909--920). Every definition, display and label of those lines is retained unchanged. Omitted from this package and not redistributed: 1--319, 323--908, 921--1062, 1099--1631. Nothing inside a retained excerpt is omitted or altered: every retained body is delivered exactly as the article's lines are written, so no omission receipt of any kind accompanies this package. Citations: the retained excerpts carry no citation command at all, so no bibliography entry of the article is transcribed and no reference is redistributed. Reference adaptation: none was needed and none was made; every reference inside the delivered unit resolves inside it, so no unresolved reference or citation is produced. Printed slips kept literal: no printed word, symbol, hypothesis or number of the counted statement or of its proof was altered. Layout and class: the article's own class is \texttt{amsart} with packages including xcolor, amsmath, amssymb, amsfonts, amscd, psfrag, graphicx, stmaryrd, epsfig, amsthm, fullpage, verbatim, listings, mathtools; the wrapper is the lane's pinned \texttt{amsart} editorial wrapper at 10pt on A4, which restates the theorem-like environments the retained text needs and adds the article's own macro definitions that the retained text needs (none), with extra wrapper packages (bm, bbm, dsfont, xspace, cancel, mathtools). The delivered numbering is the unit's own. Assets: no retained span contains a figure environment, a graphics-inclusion command, a PDF-page inclusion, a table or a TikZ picture; no image file is copied into this package, the package declares no asset dependency and no image file is redistributed with it. Licence: version arXiv:2310.15899v3 of this article is distributed under the Creative Commons Attribution 4.0 International licence, as recorded in the arXiv licence metadata for this exact version and shown on the abstract page https://arxiv.org/abs/2310.15899v3. A full rights notice is delivered at licenses/arxiv-2310.15899-CC-BY-4.0.txt. Characters: every retained excerpt is pure ASCII, so the delivered engine, XeLaTeX with its default Latin Modern text font, reports no missing character.
\begin{lemma}\label{lem:faces-in-3-face}
If $v$ is a $3$-vertex, then $m_3(v)=0$ and  $m_4(v)\leq 1$. 
\end{lemma}

\begin{remark}\label{rem:specil vertices} 
Let $v$ be a $5$-vertex.
\begin{itemize}
\item[$(a)$] If $v$ is a bad vertex, then all neighbours of $v$ are $5$-vertices, and no edge $v_iv_{i+1}$ for $i\in [4]$ is contained in two $3$-faces.
%no pair of neighbours of $v$ has a common neighbour in $G-v$.

\item[$(b)$] If $v$ is a semi-bad vertex, then $n_3(v)=0$ and $n_5(v)\geq 4$. %$v$ has at least four $5$-neighbours.
%\item[$(c)$] If $v$ is a strong  vertex, then  $v$ has no $3$-neighbour; $v$ has at least four $5$-neighbours; $v$ is incident to a $5^+$-face. 
%
%
\end{itemize}
\end{remark}

\begin{lemma}\label{lem:support-properties}
Let $v$ be a support vertex. 
\begin{itemize}
\item[$(a)$] If $n_3(v)=1$, then  $v$ is incident to a $5^+$-faces. 
\item[$(b)$] If $n_3(v)=1$ and $n_4(v)=1$ such that there exists a $5^+$-face incident to  both $v$ and its $3$-neighbour, then  $v$ is incident to two $5^+$-faces. 
\item[$(c)$] If $n_3(v)=1$ and $n_4(v)=2$, then  $v$ is incident to three $5^+$-faces. 
\item[$(d)$] If $n_3(v)=2$, then $v$ is incident to three $5^+$-faces. 
\end{itemize}
\end{lemma}

\begin{proof} 
Let $v$ be a support vertex with $3$-faces $f_1=v_1vv_2$ and $f_2=v_2vv_3$. Then $v_2$ is a bad or a semi-bad vertex, which implies that at least one of $v_1v_2$ and $v_2v_3$ is contained in two $3$-faces.

(a).  Let $n_3(v)=1$.  By Lemma \ref{lem:faces-in-3-face}, the $3$-neighbour of $v$ is either $v_4$ or $v_5$. Assume, without loss of generality,  that $v_5$ is a $3$-vertex. By contradiction, assume that all faces incident to $v$ other than $f_1,f_2$ are $4$-faces. We then say that each pair of $\{v_3,v_4\}, \{v_4,v_5\}, \{v_1,v_5\}$ has a common neighbour in $G-v$. Clearly, we have $d_2(v)\leq 15$. Thus, if we set $G'=G-v+\{v_1v_5,v_2v_5,v_4v_5\}$, then $G'$ would be proper with respect to $G$. By minimality, the graph $G'$ has a 2-distance coloring $f$ with $16$ colors. Since $v$ has at most $15$ forbidden colors, we can color $v$ with an available color, a contradiction. Hence, $v$ is incident to a $5^+$-face.

(b). Let $n_3(v)=1$ and $n_4(v)=1$. Suppose that there exists a $5^+$-face incident to  both $v$ and its $3$-neighbour.  By Lemma \ref{lem:faces-in-3-face}, the $3$-neighbour of $v$ is either $v_4$ or $v_5$. Assume, without loss of generality,  that $v_5$ is a $3$-vertex.   Denote by $f_1,f_2,\ldots,f_5$ the faces incident to $v$ in clockwise order on the plane. By our assumption, $f_4$ or $f_5$ is a $5^+$-face.
%By (a), $v$ is incident to a $5^+$-face. 
Assume to the contrary that $v$ is incident to two $4$-faces, i.e., two of $f_3,f_4,f_5$ are $4$-face. We then have $d_2(v)\leq 15$, and deduce that $f_3$ is a $4$-face.
%If $v_4$ is a $4$-vertex, then we set $G'=G-v+\{v_3v_4,v_4v_5,v_5v_1\}$, and so $G'$ is proper with respect to $G$.  Similarly as above, we get a contradiction since $d_2(v)\leq 15$. We we may assume further that $v_4$ is a $5$-vertex.  
%As $v$ is incident to two $4$-faces and one $5^+$-face, %we consider all possible cases. 
%By the assumption, either $f_4$ or $f_5$ is a $5^+$-face. 
%If $f_3$ is a $5^+$-face, then we set $G'=G-v+\{v_2v_4,v_3v_5,v_1v_5\}$.
%If $f_4$ or $f_5$ is a $5$-face, 
When we set $G'=G-v+\{v_1v_5,v_2v_5,v_4v_5\}$, $G'$ is proper with respect to $G$. Similarly as above, we get a contradiction. Hence, $v$ is incident to two $5^+$-faces.


(c).  Let $n_3(v)=1$ and $n_4(v)=2$. By Lemma \ref{lem:faces-in-3-face}, the $3$-neighbour of $v$ is either $v_4$ or $v_5$. We may suppose that $v_5$ is a $3$-vertex. Assume to the contrary that $v$ is incident to at most two $5^+$-faces. This implies that $v$ is incident to a $4$-face. Since at most one of $v_1,v_3$ is a $4$-vertex by Remark \ref{rem:specil vertices}(b), we may assume that $v_3,v_4$ are $4$-vertices due to $n_4(v)= 2$ and $n_3(v)=1$.  In such a case, we have $d_2(v)\leq 15$ and so we set $G'=G-v+\{v_3v_4,v_4v_5,v_5v_1\}$.  Clearly,  $G'$ is proper with respect to $G$.
Similarly as above, we get a contradiction. 


(d). Let $n_3(v)=2$. So, $v_4$ and $v_5$ are $3$-vertices by Lemma \ref{lem:faces-in-3-face}.  Suppose that $v$ is incident to a $4$-face.  In such a case, we have $d_2(v)\leq 15$ and so we set $G'=G-v+\{v_3v_4,v_4v_5,v_5v_1\}$.  Clearly,  $G'$ is proper with respect to $G$.
Similarly as above, we get a contradiction. 
%By minimality, the graph $G'$ has a 2-distance coloring $f$ with $16$ colors. 
%Notice that every vertex in $V(G)$, except $v$, is colored by the same color appears on the same vertex in $G'$. 
%Since $v$ has at most $15$ forbidden colors, we can color $v$ with an available color, a contradiction. 
%Thus $m_4(v)=0$, and so $v$ is incident to three $5^+$-faces. 
\end{proof}

\end{document}
